\section*{Цели и задачи работы}
\begin{itemize}
    \item Изучение зависимостей между деформациями и напряжениями при осевом сжатии пластичных, хрупких и неоднородных материалов.
    \item Определение сопротивления указанных материалов сжимающим нагрузкам.
\end{itemize}

\section*{Теоретические сведения}
Многие элементы конструкций и приборов в процессе эксплуатации подвергаются сжатию. Поведение материалов под действием сжимающих нагрузок существенно различается в зависимости от их механических свойств.

Для проведения испытаний на сжатие, как правило, используются образцы в форме коротких цилиндров или кубов. Чтобы избежать продольного изгиба (потери устойчивости) и обеспечить однородное напряженное состояние в центральной части, высота образца $h$ должна находиться в пределах $1.5d \le h \le 3d$, где $d$ --- диаметр цилиндра.

\subsection*{Пластичные материалы}
При сжатии пластичных материалов (например, низкоуглеродистая сталь, медь, алюминий) после достижения предела текучести образец не разрушается, а претерпевает значительные пластические деформации. Он укорачивается по высоте и увеличивается в поперечных размерах, принимая бочкообразную форму. Типичная диаграмма сжатия и вид образца до и после испытания представлены на \cref{fig:plastic_compression}.

\begin{figure}[H]
    \centering
    \includegraphics[width=0.5\textwidth]{figs/plastic_compression.pdf}
    \caption{Испытание образца из пластичного материала: \textit{а} --- диаграмма сжатия; \textit{б}, \textit{в} --- вид образца до и после испытания}
    \label{fig:plastic_compression}
\end{figure}

На диаграмме (см. \cref{fig:plastic_compression}, а) можно выделить два характерных участка: упругой деформации (до нагрузки $P_T$) и пластической деформации (текучести). Разрушения как такового не происходит, испытание прекращают при достижении заданной степени деформации.

\subsection*{Хрупкие материалы}
Хрупкие материалы (например, чугун, бетон, керамика) при сжатии ведут себя упруго вплоть до момента разрушения. Разрушение происходит внезапно, путем среза или раскалывания. Плоскость разрушения обычно наклонена под углом $45^\circ \dots 50^\circ$ к оси образца, что соответствует плоскостям действия максимальных касательных напряжений. Диаграмма сжатия для хрупкого материала практически линейна до максимальной нагрузки $P_\text{max}$, при которой происходит разрушение (\cref{fig:brittle_compression}).

\begin{figure}[H]
    \centering
    \includegraphics[width=0.5\textwidth]{figs/brittle_compression.pdf}
    \caption{Испытание образца из хрупкого материала: \textit{а} --- диаграмма сжатия; \textit{б}, \textit{в} --- вид образца до и после испытания}
    \label{fig:brittle_compression}
\end{figure}

\subsection*{Неоднородные (анизотропные) материалы}
Поведение неоднородных материалов, таких как древесина, зависит от направления приложения нагрузки относительно структуры материала (волокон).

\begin{itemize}
    \item \textit{Сжатие вдоль волокон.} Древесина ведет себя подобно хрупкому материалу. Разрушение происходит из-за потери устойчивости отдельных волокон и их последующего излома. На диаграмме сжатия наблюдается почти линейный рост нагрузки до пикового значения $P_\text{max}$ (\cref{fig:anisotropic_compression}, кривая 1).
    \item \textit{Сжатие поперек волокон.} В этом случае древесина проявляет свойства пластичного материала. Происходит смятие и уплотнение волокон без явного разрушения. Нагрузка $P_1$ соответствует пределу пропорциональности, после которого начинается интенсивная пластическая деформация (\cref{fig:anisotropic_compression}, кривая 2).
\end{itemize}

\begin{figure}[H]
    \centering
    \includegraphics[width=0.5\textwidth]{figs/anisotropic_compression.pdf}
    \caption{Испытание образцов из неоднородного материала (древесина): \textit{а} --- диаграмма сжатия (1 --- вдоль волокон, 2 --- поперек волокон); \textit{б}, \textit{в} --- вид образцов до и после испытания}
    \label{fig:anisotropic_compression}
\end{figure}

% Correspondence between figures in the source PDF and the generated code:
% Рис. 2.1 -> \cref{fig:plastic_compression}
% Рис. 2.2 -> \cref{fig:brittle_compression}
% Рис. 2.3 -> \cref{fig:anisotropic_compression}